\documentclass{report}
\usepackage{amsmath,amssymb}
\setlength{\parindent}{0mm}
\setlength{\parskip}{1em}
\begin{document}
\begin{center}
\rule{6in}{1pt} \
{\large Andreas Noack \\
{\bf Faster than FastPCA: High performance principal components analysis of genomics data}}

32 Vassar St \\ Rm 32-206 \\ MIT CSAIL \\ Cambridge \\ MA 02139-4307
\\
{\tt noack@mit.edu}\\
Jiahao Chen\\
Jake Bolewski\\
Alan Edelman\\
Nikolaos Patsopoulos\end{center}

Data for genome-wide association studies (GWAS) involve extremely large matrices
with small nonnegative integer entries representing deviations from a reference
genome. Principal components analysis (PCA) is traditionally used to identify
clusters reflecting subpopulation structure in genomics populations.

We present implementations of high performance principal components analysis using
Golub-Kahan-Lanczos (GKL) iterative bidiagonalization routines written in pure Julia.
We demonstrate out of core computation on data stored on external files and databases,
showing how Julia's user-extensible type system and generic functions with
multimethods (multiple dispatch) allow for flexible and convenient
definitions of new types reflecting new data structures, such as matrices stored
in HDF5 files or in an array database. Operator overloading also allows for a
natural expression of elementary operations such as matrix-vector products that
exploit detailed knowledge of underlying storage formats and layouts for improved
performance.

We compare our implementation of principal components analysis with existing
tools designed for GWAS data, such as EIGENSOFT, FlashPCA and FastPCA. We show
that in many cases, iterative bidiagonalization outperforms other methods
implemented in these existing tools, such as randomized subspace iteration.
We explore how convergence of GKL iterates and auxiliary quantities used in
practical roundoff control strategies, such as the $\omega$-recurrence coefficients
in partial reorthogonalization reflect underlying structure latent in genomics
data matrices.


\end{document}
